% Run pdfLaTeX from the source directory. Keep the lookup for the activity's year.
% Exactly ONE worked example, ONE guided example, and FIVE independent exercises.
% Counts apply to the whole activity; never repeat these blocks per family.
% Printable student activity only. Required gray headings and white criterion boxes
% are structural styles; semantic instructional highlights and legends are forbidden.
% Strip source highlighting while preserving content and ordinary formatting.
% Avoid inherited highlight-producing macros (e.g. the shaded preamble boxed macro).
\makeatletter
\def\input@path{{../../../preamble/}}
\makeatother
\input{preamble.tex}
\usepackage{needspace}
\usepackage{varwidth}

\geometry{letterpaper,margin=0.25in}
\pagestyle{empty}
\linespread{1.04}
\setlength{\parskip}{4pt}
\raggedbottom

% Exactly TWO table-style families. All helpers size to local \linewidth.
% Family 1: plain opening notation/formula reference, never a highlight legend.
% Adapt widths or use full width for the topic's actual definitions and formulas.
\newcommand{\NotationReferenceTable}[1]{%
    \par\begingroup
    \small
    \setlength{\tabcolsep}{3pt}
    \setlength{\extrarowheight}{0pt}
    \setlength{\parskip}{0pt}
    \renewcommand{\arraystretch}{1}
    \noindent\begin{tabularx}{\linewidth}{@{}>{\raggedright\arraybackslash}p{.21\linewidth}>{\raggedright\arraybackslash}p{.34\linewidth}>{\raggedright\arraybackslash}X@{}}
        \toprule
        \textbf{Symbol} & \textbf{Concept} & \textbf{Formula} \\
        \midrule
        #1
        \bottomrule
    \end{tabularx}\par
    \endgroup
}
% Family 2: all instructional/problem tables, in the printable activity.
% Width variants below are the SAME ruled family; adapt column count if needed.
% Body rows end with \\ \hline. Keep group/case headings short or wrap them.
\newcommand{\ClassroomRuledTable}[4]{%
    \par\begingroup
    \small
    \setlength{\tabcolsep}{4pt}
    \setlength{\extrarowheight}{0pt}
    \setlength{\parskip}{0pt}
    \renewcommand{\arraystretch}{1.16}
    \noindent\begin{tabularx}{\linewidth}{|>{\raggedright\arraybackslash}p{#1}|>{\raggedright\arraybackslash}p{#2}|>{\raggedright\arraybackslash}X|}
        \hline
        #3 \\ \hline
        #4
    \end{tabularx}\par
    \endgroup
}
\newcommand{\InformationTable}[1]{%
    \ClassroomRuledTable{.14\linewidth}{.32\linewidth}
        {\textbf{Symbol} & \textbf{Concept} & \textbf{In context...}}{#1}%
}
\NewDocumentCommand{\CalculationTable}{O{.10\linewidth} O{.20\linewidth} +m}{%
    \ClassroomRuledTable{#1}{#2}
        {\textbf{Symbol} & \textbf{Concept} & \textbf{Calculation, result, and meaning in context}}{#3}%
}
\newcommand{\TableGroup}[1]{\multicolumn{3}{|l|}{\textbf{#1}}\\ \hline}
% Structural styles for the printable activity only; no semantic highlighting.
\newcommand{\stageheading}[1]{%
    \par\Needspace{4\baselineskip}\addvspace{4pt}%
    {\centering\large\bfseries
     \tikz[baseline=(title.base)]{%
         \node[fill=black!10,text=black,rounded corners=2mm,
             inner sep=3pt,outer sep=0pt] (title) {\strut #1};}\par}%
    \nopagebreak[4]\smallskip
}
\newcommand{\criterionbox}[2]{%
    \begin{tcolorbox}[colback=white,colframe=black!45,boxrule=.4pt,
        arc=2mm,boxsep=0pt,left=6pt,right=6pt,top=5pt,bottom=5pt,
        before skip=4pt,after skip=4pt,width=\linewidth]
        \textbf{#1} #2
    \end{tcolorbox}%
}
\newcommand{\problemtitle}[3]{%
    \par\Needspace{4\baselineskip}\addvspace{0.4\baselineskip}%
    \noindent
    {\begin{varwidth}[t]{0.75\linewidth}%
         \bfseries\strut #1. #2%
     \end{varwidth}}%
    \hfill\textbf{#3}\par\nopagebreak[4]
}
\newcommand{\taskheading}[1]{%
    \par\Needspace{5\baselineskip}\addvspace{4pt}%
    \noindent\tikz[baseline=(title.base)]{%
        \node[fill=black!5,text=black,rounded corners=2mm,
            inner sep=3pt,outer sep=0pt] (title) {%
            \begin{varwidth}{\dimexpr\linewidth-6pt\relax}%
                \raggedright\bfseries\strut #1\strut
            \end{varwidth}};}\par\nopagebreak[4]\smallskip
}
\newcommand{\task}[2]{\taskheading{#1}#2}
\newcommand{\responseline}{\par\noindent\rule{0pt}{1.55em}\hrulefill\par}
\newcommand{\workspace}[1]{\par\noindent\rule{0pt}{#1}\par\nopagebreak[4]}


% Local graphs use only the packages in the year-local preamble.
\newcommand{\EVGraph}[4]{%
\begin{tikzpicture}\begin{axis}[
 width=.88\linewidth,height=5.4cm,xmin=0,xmax=1,
 ymin={min(#2,#3)-.2*abs(#3-#2)},ymax={max(#2,#3)+.2*abs(#3-#2)},
 axis lines=left,xlabel={$p=P(E)$},ylabel={Expected net profit (USD)},
 title={$\mathbb E[Y_{#1}]$},xtick={0,1},ytick={#2,0,#3},
 extra x ticks={#4},extra x tick labels={$#4$},
 extra x tick style={tick label style={font=\small}},grid=major]
 \addplot[black,thick,domain=0:1,samples=2]{#2+(#3-#2)*x};
 \addplot[black,dashed,domain=0:1,samples=2]{0};
 \addplot[black,only marks,mark=*] coordinates {(#4,0)};
\end{axis}\end{tikzpicture}}
\newcommand{\BlankGraph}[1]{%
\begin{tikzpicture}[x=6.1cm,y=1cm]
 \draw[step=.61cm,black!20,very thin] (0,-1.4) grid (1,1.4);
 \draw[->] (0,-1.4)--(0,1.7) node[above,font=\small] {$\mathbb E[Y_{#1}]$ (USD)};
 \draw[->] (0,0)--(1.08,0) node[right,font=\small] {$p$};
 \node[below left] at (0,0) {$0$}; \node[below] at (1,0) {$1$};
\end{tikzpicture}}

% Short corresponding blocks can break naturally between paired procedures.
\NewDocumentCommand{\ColumnPair}{+m +m}{%
\par\noindent\begin{minipage}[t]{.49\linewidth}\vspace{0pt}\raggedright
#1
\end{minipage}\hfill
\begin{minipage}[t]{.49\linewidth}\vspace{0pt}\raggedright
#2
\end{minipage}\par}
\newcommand{\calcblank}{\underline{\hspace{1.15cm}}}
\newcommand{\shortblank}{\underline{\hspace{.7cm}}}
\newcommand{\PayoffSummary}[1]{%
\par\begingroup\small\setlength{\tabcolsep}{4pt}
\setlength{\parskip}{0pt}\renewcommand{\arraystretch}{1.16}
\begin{tabularx}{\linewidth}{|>{\raggedright\arraybackslash}p{.44\linewidth}|>{\raggedright\arraybackslash}X|}\hline
\textbf{Payoff summary} & \textbf{Value}\\\hline
#1\end{tabularx}\par\endgroup}
\begin{document}

{\centering {\Large\bfseries Bernoulli Models and Individual Break-Even Analysis\par}
Global Economics \textbar{} Grade 10 \textbar{} Term 2 \textbar{} Lesson 1\par}
\criterionbox{C3:}{Models binary economic outcomes with Bernoulli random variables, interpreting success and failure in contexts such as default, gain, loss, or claim occurrence.}

\stageheading{Equations and basic procedure}
One observation in its stated period has two mutually exclusive, exhaustive states.
The event label ``success'' means that $E$ occurs; it need not mean a monetary gain.
All data are hypothetical. All monetary payoffs are net gains or losses after the stated costs. Each $Y_i$ has one strictly positive and one strictly negative payoff.
\NotationReferenceTable{
$E, E^c$ & Event and complement & $E^c$: all outcomes outside $E$.\\
$X,p$ & Unitless Bernoulli indicator & $X=1$ on $E$, $X=0$ on $E^c$; $p=P(E)\in[0,1]$.\\
$1-p$ & Complementary mass & $P(X=0)=1-p$; $P(X=1)=p$.\\
$S_X,S_{Y_i}$ & Positive-probability supports & $S_X=\{x\in\mathbb R:P(X=x)>0\}$; $S_{Y_i}=\{y\in\mathbb R:P(Y_i=y)>0\}$.\\
$f_i,Y_i$ & State-to-money rule (USD) & $f_i:\{0,1\}\to\mathbb R$; $Y_i=f_i(X)$.\\
$f_i(0),f_i(1)$ & Fixed state payoffs (USD) & $Y_i=f_i(0)+[f_i(1)-f_i(0)]X$; $f_i(0)f_i(1)<0$.\\
$m_Z(z)$ & Probability mass function & $m_Z(z)=P(Z=z)$; zero outside the stated values.\\
$F_Z(t)$ & Cumulative distribution function & $F_Z(t)=P(Z\leq t)$; includes mass at the threshold.\\
$\mathbb E[Y_i]$ & Expected net profit (USD) & $p f_i(1)+(1-p)f_i(0)=f_i(0)+[f_i(1)-f_i(0)]p$.\\
$p_i^*$ & Individual break-even probability & $p_i^*=-f_i(0)/[f_i(1)-f_i(0)]\in(0,1)$.\\
}
\ColumnPair{
\textbf{Indicator distribution (valid also at $p=0,1$):}
\[
\begin{gathered}
m_X(x)\\
=
\end{gathered}
\begin{cases}1-p&x=0\\p&x=1\\0&x\notin\{0,1\}\end{cases}
\enspace
\begin{gathered}
F_X(t)\\
=
\end{gathered}
\begin{cases}0&t<0\\1-p&0\leq t<1\\1&t\geq1\end{cases}
\]
}{
\textbf{Monetary distribution:} write $a=\min(f_i(0),f_i(1))$, $b=\max(f_i(0),f_i(1))$ and $q=P(Y_i=a)$.
\[
\begin{gathered}
m_{Y_i}(y)\\
=
\end{gathered}
\begin{cases}1-p&y=f_i(0)\\p&y=f_i(1)\\0&\text{otherwise}\end{cases}
\enspace
\begin{gathered}
F_{Y_i}(t)\\
=
\end{gathered}
\begin{cases}0&t<a\\q&a\leq t<b\\1&t\geq b\end{cases}
\]
}
For $0<p<1$, $S_X=\{0,1\}$ and $S_{Y_i}=\{f_i(0),f_i(1)\}$.

\stageheading{General tasks}\vspace{-3mm}
\noindent\begin{minipage}[t]{.49\linewidth}\vspace{0pt}
\task{Task 1 --- Binary model and distributions.}{}
\begin{enumerate}[label=\arabic*.,itemsep=2pt,topsep=2pt,leftmargin=1.5em]
\item State one observation, the states $E,E^c$, the unitless $X$, and monetary rules $Y_i=f_i(X)$ with state payoffs (USD).
\item State support sets, complete PMFs, and numerically ordered piecewise CDFs for $X$ and each $Y_i$.
\item Each loss event $\{Y_i<0\}$ in $X$, its probability function and contextual meaning.
\end{enumerate}
\end{minipage}\hfill
\begin{minipage}[t]{.49\linewidth}\vspace{0pt}
\task{Task 2 --- Individual expected values.}{}
\begin{enumerate}[label=\arabic*.,itemsep=2pt,topsep=2pt,leftmargin=1.5em]
\item Calculate each expectation at the initial probability (USD) and interpret it per period.
\item For fixed payoffs and varying $p\in[0,1]$, derive each expectation formula, exact break-even root, and positive/negative intervals; summarize in tables.
\item Separate labelled graphs: $p\in[0,1]$, USD scale, endpoints, expectation and zero lines, exact roots.
\end{enumerate}
\end{minipage}\par


\stageheading{Worked example}\vspace{-3mm}
\problemtitle{1}{Campus concert food stand or indoor games}{}
A student club compares two hypothetical fundraising choices for one weekend: a food
stand at a campus outdoor concert or an indoor games afternoon. The concert either goes
ahead or is cancelled; initially, it goes ahead on three out of every eight weekends.
After food, room, equipment, and publicity costs as applicable, the food stand has a net
gain of USD 105 if the concert goes ahead and a net loss of USD 45 if it is cancelled.
The games afternoon has a net loss of USD 30 if the concert goes ahead, because it draws
potential players away, and a net gain of USD 90 if it is cancelled, because students
seek an indoor alternative. Only one option would be chosen; both use the same concert outcome.

\task{Task 1 --- Establish the binary model and monetary transformations.}{}
\textbf{\underline{1. Binary model and monetary transformations.}}\nopagebreak[4]
\textbf{Shared model.} One observation is the club's net outcome over one weekend for the chosen option.
Let $E$ be the concert going ahead and $E^c$ its cancellation.
Set $X=1$ on $E$, $X=0$ on $E^c$. Then $X\sim\operatorname{Bernoulli}(p)$ , $p=P(E)=3/8$.
\vspace{-3mm}
\ColumnPair{
\textbf{$Y_1$: food stand net profit (USD)}
\vspace{-2mm}
\[
\begin{aligned}
&f_1(0)=-45<0 \qquad f_1(1)=105>0\\
Y_1&=f_1(0)+[f_1(1)-f_1(0)]X\\
&=-45+[105-(-45)]X = -45+150X
\end{aligned}
\]
}{
\textbf{$Y_2$: games afternoon net profit (USD)}
\vspace{-2mm}
\[
\begin{aligned}
&f_2(0)=90>0 \qquad f_2(1)=-30<0\\
Y_2&=f_2(0)+[f_2(1)-f_2(0)]X\\
&=90+[(-30)-90]X = 90-120X
\end{aligned}
\]
}
\vspace{3mm}
\textbf{\underline{2. Supports, PMFs, and ordered CDFs}}\nopagebreak[4]\vspace{-6mm}
\[
\text{For } 0<p<1 \quad:\quad S_X=\{x\in\mathbb R:P(X=x)>0\}=\{0,1\} \qquad
\begin{gathered}
m_X(x)\\
=
\end{gathered}
\begin{cases}1-p&x=0\\p&x=1\\0&\text{otherwise}\end{cases}
\enspace
\begin{gathered}
F_X(t)\\
=
\end{gathered}
\begin{cases}0&t<0\\1-p&0\leq t<1\\1&t\geq1\end{cases}
\]

\vspace{-5mm}
 $ p=0 \rightarrow$  $S_X=\{0\}$ \qquad  $ p=1 \rightarrow$  $S_X=\{1\}$

\ColumnPair{
\textbf{Support} \vspace{-3mm}
\[
\begin{aligned}
S_{Y_1}&=\{y\in\mathbb R:P(Y_1=y)>0\} =\{-45,105\}
\end{aligned}
\]
}{
\textbf{Support} \vspace{-3mm}
\[
\begin{aligned}
S_{Y_2}&=\{y\in\mathbb R:P(Y_2=y)>0\} =\{-30,90\}
\end{aligned}
\]
}
\ColumnPair{
\textbf{$Y_1$ --- Complete PMF and ordered CDF.}
\[
\begin{gathered}
m_{Y_1}(y)\\
=
\end{gathered}
\begin{cases}1-p&y=-45\\p&y=105\\0&\text{otherwise}\end{cases}
\enspace
\begin{gathered}
F_{Y_1}(t)\\
=
\end{gathered}
\begin{cases}0&t<-45\\1-p&-45\leq t<105\\1&t\geq105\end{cases}
\]
}{
\textbf{$Y_2$ --- Complete PMF and ordered CDF.}
\[
\begin{gathered}
m_{Y_2}(y)\\
=
\end{gathered}
\begin{cases}p&y=-30\\1-p&y=90\\0&\text{otherwise}\end{cases}
\enspace
\begin{gathered}
F_{Y_2}(t)\\
=
\end{gathered}
\begin{cases}0&t<-30\\p&-30\leq t<90\\1&t\geq90\end{cases}
\]
}
\textbf{\underline{3. Initial probabilities and loss events.}}\nopagebreak[4]
Initially $P(X=0)=5/8$, $P(X=1)=3/8$.
\ColumnPair{
\textbf{$Y_1$ --- Initial probabilities and loss.}
\[
\begin{gathered}
m_{Y_1}(-45)=\tfrac58,\quad m_{Y_1}(105)=\tfrac38\\
\{Y_1<0\}=\{X=0\}=E^c,\quad P(Y_1<0)=1-p
\end{gathered}
\]
Five out of eight weekends produce a food stand loss.
\par
}{
\textbf{$Y_2$ --- Initial probabilities and loss.}
\[
\begin{gathered}
m_{Y_2}(-30)=\tfrac38,\quad m_{Y_2}(90)=\tfrac58\\
\{Y_2<0\}=\{X=1\}=E,\quad P(Y_2<0)=p
\end{gathered}
\]
Three out of eight weekends produce a games afternoon loss.
\par
}

\Needspace{6\baselineskip}
\task{Task 2 --- Individual expected values.}{}
\textbf{\underline{1. Initial expected values}} ($p=3/8$)\nopagebreak[4] 
\ColumnPair{
\textbf{$Y_1$: initial expectation (USD).}
\[
\begin{aligned}
\mathbb E[Y_1]\big|_{p=3/8}&=105(\tfrac38)-45(\tfrac58)\\
&=\tfrac{315-225}{8}=\tfrac{45}{4}=11.25\ \text{USD}
\end{aligned}
\]
Food stand: expected gain of USD 11.25 per weekend.

}{
\textbf{$Y_2$: initial expectation (USD).}
\[
\begin{aligned}
\mathbb E[Y_2]\big|_{p=3/8}&=-30(\tfrac38)+90(\tfrac58)\\
&=\tfrac{-90+450}{8}=\tfrac{360}{8}=45\ \text{USD}
\end{aligned}
\]
Games afternoon: expected gain of USD 45 per weekend.

}
\textbf{\underline{2. Expectation formulas and sign intervals.}} Keep payoffs fixed; vary $p\in[0,1]$.\nopagebreak[4]
\ColumnPair{
\textbf{$Y_1$: expectation and break-even (USD).}
\[
\begin{aligned}
\mathbb E[Y_1]&=p f_1(1)+(1-p)f_1(0)\\
&=105p-45(1-p)=-45+150p
\end{aligned}
\]
\[
\begin{aligned}
\mathbb E[Y_1]=0&\iff-45+150p=0\\
&\iff 150p=45\iff p=\tfrac{45}{150}=\tfrac3{10}
\end{aligned}
\]
}{
\textbf{$Y_2$: expectation and break-even (USD).}
\[
\begin{aligned}
\mathbb E[Y_2]&=p f_2(1)+(1-p)f_2(0)\\
&=-30p+90(1-p)=90-120p
\end{aligned}
\]
\[
\begin{aligned}
\mathbb E[Y_2]=0&\iff90-120p=0\\
&\iff -120p=-90\iff p=\tfrac{-90}{-120}=\tfrac34
\end{aligned}
\]
}
\ColumnPair{
\textbf{$Y_1$: sign intervals.}
\[
\begin{aligned}
\mathbb E[Y_1]>0&\iff150p>45\\
&\iff p>\tfrac3{10}\iff p\in(\tfrac3{10},1]\\[3pt]
\mathbb E[Y_1]<0&\iff150p<45\\
&\iff p<\tfrac3{10}\iff p\in[0,\tfrac3{10})
\end{aligned}
\]
Food stand: expected gain above $3/10$, loss below it, and break-even at $3/10$.

}{
\textbf{$Y_2$: sign intervals.}
\[
\begin{aligned}
\mathbb E[Y_2]>0&\iff-120p>-90\\
&\iff p<\tfrac34\iff p\in[0,\tfrac34)\\[3pt]
\mathbb E[Y_2]<0&\iff-120p<-90\\
&\iff p>\tfrac34\iff p\in(\tfrac34,1]
\end{aligned}
\]
Dividing by $-120$ reverses both inequalities.
Games afternoon: expected gain below $3/4$, loss above it, and break-even at $3/4$.

}

\textbf{\underline{3. Graphs.}}
\ColumnPair{
\textbf{$Y_1$: graph.}\par\centering
\EVGraph{1}{-45}{105}{3/10}
}{
\textbf{$Y_2$: graph.}\par\centering
\EVGraph{2}{90}{-30}{3/4}
}
\textbf{Checks.} Endpoint expectations equal the state payoffs:
$\mathbb E[Y_1](0)=-45$, $\mathbb E[Y_1](1)=105$;
$\mathbb E[Y_2](0)=90$, $\mathbb E[Y_2](1)=-30$ (USD).
At the roots, $-45+150(3/10)=0$ and $90-120(3/4)=0$.
Each weekend yields a state payoff, so a positive expectation does not guarantee profit;
at break-even, the long-run average is zero despite nonzero realized payoffs.

\stageheading{Guided Example}\vspace{-6mm}
\problemtitle{2}{Graduation photo booth options}{Scaffolded}
A student photographer compares two hypothetical photo booth options for one graduation
weekend: a simple backdrop or a decorated set. Advance bookings either meet the advertised
target or fall short; initially, the target is met on two out of every seven weekends.
After equipment hire, decorations, printing, and travel costs as applicable, the simple
backdrop has a net gain of USD 72 if the target is met and a net loss of USD 24 otherwise.
The decorated set has a net gain of USD 126 if the target is met and a net loss of USD 54
otherwise: it attracts higher-priced bookings but costs more to prepare. Only one option
would be chosen; both are modelled using the same booking-target outcome.

\task{Task 1 --- Establish the binary model and monetary transformations.}{}
\textbf{Shared model.} One observation is the photographer's net outcome over one graduation weekend.
Let $E$ be meeting the booking target and $E^c$ falling short; exactly one occurs.
Set $X=1$ on $E$, $X=0$ on $E^c$. Then $X\sim\operatorname{Bernoulli}(p)$ is unitless,
$p=P(E)$, initially $p=2/7$, and both options use the same $X$; no independence is assumed.
Complete the numerical entries throughout the scaffold.
\[
S_X=\{x\in\mathbb R:P(X=x)>0\}=\{0,1\}\quad(0<p<1),\quad
\begin{gathered}
m_X(x)\\
=
\end{gathered}
\begin{cases}1-p&x=0\\p&x=1\\0&\text{otherwise}\end{cases}
\enspace
\begin{gathered}
F_X(t)\\
=
\end{gathered}
\begin{cases}0&t<0\\1-p&0\leq t<1\\1&t\geq1\end{cases}
\]
Initially $P(X=1)=2/7$, $P(X=0)=1-2/7=\calcblank$;
the CDF plateau is $\calcblank$.
At $p=0$, $S_X=\{\shortblank\}$; at $p=1$, $S_X=\{\shortblank\}$.

\ColumnPair{
\textbf{$Y_1$: simple-backdrop net profit (USD).}\par
\textbf{State payoffs and affine rule.}
\[
\begin{aligned}
f_1(0)&=-24<0,& f_1(1)&=72>0\\
Y_1&=f_1(0)+[f_1(1)-f_1(0)]X\\
&=-24+[72-(-24)]X\\
&=-24+\calcblank X
\end{aligned}
\]
\textbf{Support} (USD; $0<p<1$).
\[
\begin{aligned}
S_{Y_1}&=\{y\in\mathbb R:P(Y_1=y)>0\}\\
&=\{\calcblank,\calcblank\}
\end{aligned}
\]
}{
\textbf{$Y_2$: decorated-set net profit (USD).}\par
\textbf{State payoffs and affine rule.}
\[
\begin{aligned}
f_2(0)&=-54<0,& f_2(1)&=126>0\\
Y_2&=f_2(0)+[f_2(1)-f_2(0)]X\\
&=-54+[126-(-54)]X\\
&=-54+\calcblank X
\end{aligned}
\]
\textbf{Support} (USD; $0<p<1$).
\[
\begin{aligned}
S_{Y_2}&=\{y\in\mathbb R:P(Y_2=y)>0\}\\
&=\{\calcblank,\calcblank\}
\end{aligned}
\]
}
\ColumnPair{
\textbf{$Y_1$ --- Complete PMF and ordered CDF.}
\[
\begin{aligned}
&\begin{gathered}
m_{Y_1}(y)\\
=
\end{gathered}
\begin{cases}1-p&y=\calcblank\\[5pt]p&y=\calcblank\\[5pt]0&\text{otherwise}\end{cases}\\[5pt]
&\begin{gathered}
F_{Y_1}(t)\\
=
\end{gathered}
\begin{cases}0&t<\calcblank\\[5pt]1-p&\shortblank\leq t<\shortblank\\[5pt]1&t\geq\calcblank\end{cases}
\end{aligned}
\]
}{
\textbf{$Y_2$ --- Complete PMF and ordered CDF.}
\[
\begin{aligned}
&\begin{gathered}
m_{Y_2}(y)\\
=
\end{gathered}
\begin{cases}1-p&y=\calcblank\\[5pt]p&y=\calcblank\\[5pt]0&\text{otherwise}\end{cases}\\[5pt]
&\begin{gathered}
F_{Y_2}(t)\\
=
\end{gathered}
\begin{cases}0&t<\calcblank\\[5pt]1-p&\shortblank\leq t<\shortblank\\[5pt]1&t\geq\calcblank\end{cases}
\end{aligned}
\]
}
\ColumnPair{
\textbf{$Y_1$ --- Initial probabilities and loss.}
\[
\begin{gathered}
m_{Y_1}(-24)\\
=
\end{gathered}
1-\tfrac27=\calcblank,
\quad
\begin{gathered}
m_{Y_1}(72)\\
=
\end{gathered}
\calcblank
\]
$\{Y_1<0\}=\{X=\shortblank\}=E^c$.\par
$P(Y_1<0)=1-p=1-2/7=\calcblank$ initially.\par
In the long run, $\shortblank$ out of seven weekends produce a simple-backdrop loss.
\par\textbf{Endpoint supports} (USD).
\[
p=0:\ S_{Y_1}=\{\calcblank\};\quad p=1:\ S_{Y_1}=\{\calcblank\}
\]
}{
\textbf{$Y_2$ --- Initial probabilities and loss.}
\[
\begin{gathered}
m_{Y_2}(-54)\\
=
\end{gathered}
1-\tfrac27=\calcblank,
\quad
\begin{gathered}
m_{Y_2}(126)\\
=
\end{gathered}
\calcblank
\]
$\{Y_2<0\}=\{X=\shortblank\}=E^c$.\par
$P(Y_2<0)=1-p=1-2/7=\calcblank$ initially.\par
In the long run, $\shortblank$ out of seven weekends produce a decorated-set loss.
\par\textbf{Endpoint supports} (USD).
\[
p=0:\ S_{Y_2}=\{\calcblank\};\quad p=1:\ S_{Y_2}=\{\calcblank\}
\]
}
\task{Task 2 --- Individual expected values.}{}
\textbf{\underline{1. Initial expected values}} ($p=2/7$).\nopagebreak[4]
\ColumnPair{
\textbf{$Y_1$: initial expectation (USD).}
\[
\begin{aligned}
\mathbb E[Y_1]\big|_{p=2/7}&=72(\tfrac27)-24(1-\tfrac27)\\
&=\frac{144-24(\shortblank)}7\\[5pt]
&=\calcblank\ \text{USD}
\end{aligned}
\]
Expected \underline{\hspace{1.4cm}} (gain/loss): USD $\calcblank$ per weekend.

}{
\textbf{$Y_2$: initial expectation (USD).}
\[
\begin{aligned}
\mathbb E[Y_2]\big|_{p=2/7}&=126(\tfrac27)-54(1-\tfrac27)\\
&=\frac{252-54(\shortblank)}7\\[5pt]
&=\calcblank\ \text{USD}
\end{aligned}
\]
Expected \underline{\hspace{1.4cm}} (gain/loss): USD $\calcblank$ per weekend.

}
\textbf{\underline{2. Expectation formulas and sign intervals.}} Keep payoffs fixed; vary $p\in[0,1]$.\nopagebreak[4]
\ColumnPair{
\textbf{$Y_1$: expectation and break-even (USD).}
\[
\begin{aligned}
\mathbb E[Y_1]&=p f_1(1)+(1-p)f_1(0)\\
&=72p-24(1-p)=-24+\calcblank p
\end{aligned}
\]
\[
\begin{aligned}
\mathbb E[Y_1]=0&\iff-24+[72-(-24)]p=0\\
&\iff\calcblank p=24\\[5pt]
&\iff p=\frac{24}{\shortblank}=\calcblank
\end{aligned}
\]
}{
\textbf{$Y_2$: expectation and break-even (USD).}
\[
\begin{aligned}
\mathbb E[Y_2]&=p f_2(1)+(1-p)f_2(0)\\
&=126p-54(1-p)=-54+\calcblank p
\end{aligned}
\]
\[
\begin{aligned}
\mathbb E[Y_2]=0&\iff-54+[126-(-54)]p=0\\
&\iff\calcblank p=54\\[5pt]
&\iff p=\frac{54}{\shortblank}=\calcblank
\end{aligned}
\]
}
\ColumnPair{
\textbf{$Y_1$: sign intervals.}
\[
\begin{aligned}
\mathbb E[Y_1]>0&\iff[72-(-24)]p>24\\
&\iff p>\calcblank\iff p\in(\calcblank,1]\\[5pt]
\mathbb E[Y_1]<0&\iff[72-(-24)]p<24\\
&\iff p<\calcblank\iff p\in[0,\calcblank)
\end{aligned}
\]
Simple backdrop: expected gain above $\calcblank$, loss below $\calcblank$, and break-even at $\calcblank$.

}{
\textbf{$Y_2$: sign intervals.}
\[
\begin{aligned}
\mathbb E[Y_2]>0&\iff[126-(-54)]p>54\\
&\iff p>\calcblank\iff p\in(\calcblank,1]\\[5pt]
\mathbb E[Y_2]<0&\iff[126-(-54)]p<54\\
&\iff p<\calcblank\iff p\in[0,\calcblank)
\end{aligned}
\]
Decorated set: expected gain above $\calcblank$, loss below $\calcblank$, and break-even at $\calcblank$.

}
\ColumnPair{
\PayoffSummary{$f_1(0),f_1(1)$ (USD)&$-24,72$\\\hline
$\mathbb E[Y_1]$ (USD)&$-24+\shortblank p$\\[5pt]\hline
Exact $p_1^*$&$\calcblank$\\[5pt]\hline
Initial value (USD)&$\calcblank$\\[5pt]\hline}
}{
\PayoffSummary{$f_2(0),f_2(1)$ (USD)&$-54,126$\\\hline
$\mathbb E[Y_2]$ (USD)&$-54+\shortblank p$\\[5pt]\hline
Exact $p_2^*$&$\calcblank$\\[5pt]\hline
Initial value (USD)&$\calcblank$\\[5pt]\hline}
}
\textbf{\underline{3. Graphs.}} Label the axes and each USD scale; join the endpoint payoffs, show the zero line, and mark the exact root.
\ColumnPair{
\textbf{$Y_1$: graph.}\par
\par\centering\BlankGraph{1}\par Exact zero: $p_1^*=\calcblank$.
}{
\textbf{$Y_2$: graph.}\par
\par\centering\BlankGraph{2}\par Exact zero: $p_2^*=\calcblank$.
}


\stageheading{Independent Practice}\vspace{-3mm}
\problemtitle{3}{Campus portraits or an online CV workshop}{Core: service}
A recent graduate compares two hypothetical freelance activities for one day: campus
portrait sessions or an online CV workshop. The university's orientation day either goes
ahead or is cancelled; initially, it goes ahead on three out of every five days considered.
After equipment, software, and advertising costs as applicable, portraits have a net gain
of USD 84 if orientation goes ahead and a net loss of USD 36 if it is cancelled. The CV
workshop has a net loss of USD 48 if orientation goes ahead, since potential participants
are busy on campus, and a net gain of USD 72 if it is cancelled, since they are free to
attend online. Only one activity would be chosen; both use the same orientation outcome.

\problemtitle{4}{Meal preparation in a shared kitchen}{Core: sales}
Housemates compare two hypothetical plans for a small meal-preparation side business over
one month: disposable food boxes or reusable containers. Monthly subscriptions either
meet their advertised target or fall short; initially, the target is met in two out of
every five months. After ingredients, packaging, kitchen use, and cleaning costs as
applicable, the disposable-box plan has a net gain of USD 98 if the target is met and a
net loss of USD 42 otherwise. The reusable-container plan has a net gain of USD 70 if the
target is met and a net loss of USD 70 otherwise, reflecting its container and cleaning
costs. Only one plan would be chosen; both use the same subscription-target outcome.

\Needspace{12\baselineskip}
\problemtitle{5}{Freelance video editing equipment}{Core: capacity}
A freelance video editor compares two hypothetical equipment choices over one quarter:
a personal laptop or a workstation rented only when a large project is confirmed.
A large project either arrives or does not arrive; initially, one arrives in four out
of every nine quarters. After software, promotion, and any equipment rental costs, the
laptop option has a net gain of USD 150 with the project and a net loss of USD 75 without
it. The workstation option has a net gain of USD 375 with the project, because faster
editing allows more paid deliverables, and a net loss of USD 75 without it, when no
workstation is rented. Only one option would be chosen; both use the same project outcome.

\problemtitle{6}{First-job event shifts}{Extension: precision}
A young adult compares two hypothetical job offers for one paid shift during a week:
catering work or ticket-desk work at the same event. The event either is cancelled or
goes ahead; initially, cancellation occurs in five out of every twelve weeks. After
transport bookings and required preparation costs, the catering offer has a net loss of
USD 27.50 if the event is cancelled and a net gain of USD 82.50 if it goes ahead. The
ticket-desk offer has a net loss of USD 45 if it is cancelled and a net gain of USD 67.50
if it goes ahead. Cancellation means no shift pay, but preparation and booked transport
still cost money. Only one offer would be accepted; both use the same weekly event outcome.
The decimal amounts are exact.

\problemtitle{7}{Online tutoring or an offline study workshop}{Extension: relabelling}
A student tutor compares two hypothetical choices for a scheduled revision weekend within
one month: online tutoring from a shared apartment or an offline study workshop at a
nearby library. The apartment building's internet either is interrupted during that
weekend or remains uninterrupted; initially, an interruption occurs in seven out of every
twenty months. After preparation, advertising, internet, and room-booking costs as
applicable, online tutoring has a net loss of USD 64 with interruption and a net gain of
USD 96 with uninterrupted service. The offline workshop has a net gain of USD 112 with
interruption, as affected students seek help at the library, and a net loss of USD 28
with uninterrupted service, when too few attend. Only one option would be chosen; both
use the same monthly internet outcome. For each option, also express net profit as a
constant plus a coefficient times an indicator that equals one for uninterrupted service
and zero for interruption. Express its expected net profit in terms of the probability of
uninterrupted service, which is one minus the interruption probability; find the exact
break-even probability of uninterrupted service and calculate its initial expected net profit
under this relabelling. Use the interruption probability on the horizontal axis of both
expected-value graphs.

\end{document}
